% !TEX root = ../algebra_02.tex
\section{Разрешимость конечной $p$-группы}

\subsection*{Определение разрешимой группы}
Напомним: группа $G$ называется \textit{разрешимой}, если существует нормальный ряд
\[
G = G_n \supset G_{n-1} \supset \cdots \supset G_1 \supset G_0 = \{e\},
\]
такой что $G_{i-1} \trianglelefteq G_i$ и $G_i/G_{i-1}$ абелева для всех $i = 1, \ldots, n$.

\begin{Theorem}{}{psolv}
	Любая конечная $p$-группа разрешима.
\end{Theorem}

\begin{proof}
	Пусть $|G| = p^n$. Доказательство проведем индукцией по $n$.

	\textit{База:} $n = 1$.
	Тогда $G \cong \mathbb{Z}/p\mathbb{Z}$ — абелева, следовательно, разрешима.

	\textit{Шаг индукции.} Предположим, что все $p$-группы порядка $p^m$, $m < n$, разрешимы.
	Так как $|G| = p^n$, по теореме о центре $p$-групп $Z(G) \neq \{e\}$.
	Положим $G' = G/Z(G)$; тогда $|G'| = p^m$ при некотором $m < n$.
	По гипотезе индукции $G'$ разрешима:
	\[
	G' = G'_l \supset G'_{l-1} \supset \cdots \supset G'_0 = \{e\}, \quad G'_k/G'_{k-1} \text{ — абелева.}
	\] 
	Пусть $\pi \colon G \to G'$ — каноническая проекция, $H = Z(G) = \ker\pi$.
	Определим $G_k = \pi^{-1}(G'_k)$ для $k = 0, 1, \ldots, l$.
	Тогда $G'_k = G_k/H$, и по второй теореме об изоморфизме:
	\[
	G'_k/G'_{k-1} = (G_k/H)\big/(G_{k-1}/H) \cong G_k/G_{k-1} \quad \text{— абелева.}
	\] 
	Ряд
	\[
	G = G_l \supset G_{l-1} \supset \cdots \supset G_1 \supset G_0 = H \supset \{e\}
	\]
	имеет абелевы факторы (поскольку $H = Z(G)$ абелева), поэтому $G$ разрешима.
\end{proof}
